
\subsection{RQ7: What does \toolname\ cost in time, tokens, API calls, and memory?}
\label{subsec:rq7-cost}

Costs are measured on the 400 goal-independent programs, with all methods
under Pass@1. \rev{Fig.~\ref{fig:rq7-cost} compares every attempt, including
failures, on \texttt{gpt-5.4-mini}; the original \autospec\ time, call count,
and validity labels are retained.} Peak RSS stays below \(\sim 500\)~MB.

\begin{figure*}[t]
\centering
\begin{tikzpicture}[x=1cm,y=0.50cm,font=\footnotesize,every node/.style={text=figgraytext}]
  \def\coststack#1#2#3#4#5#6#7{%
    \path[fill=#5] (#1,{#2-0.18}) rectangle (#4,{#2+0.18});
    \path[pattern=north east lines,pattern color=figgraytext!65]
      (#3,{#2-0.18}) rectangle (#4,{#2+0.18});
    \draw[draw=#6,line width=0.45pt] (#1,{#2-0.18}) rectangle (#4,{#2+0.18});
    \draw[draw=#6,line width=0.35pt] (#3,{#2-0.18}) -- (#3,{#2+0.18});
    \node[anchor=west] at ({#4+0.05},#2) {\rev{#7}};
  }

  \path[fill=figcontractfill] (4.55,5.88) rectangle (4.85,6.18);
  \draw[draw=figcontractstroke] (4.55,5.88) rectangle (4.85,6.18);
  \node[anchor=west] at (4.91,6.03) {\rev{Successful runs}};
  \path[fill=figcontractfill] (7.30,5.88) rectangle (7.60,6.18);
  \path[pattern=north east lines,pattern color=figgraytext!65]
    (7.30,5.88) rectangle (7.60,6.18);
  \draw[draw=figcontractstroke] (7.30,5.88) rectangle (7.60,6.18);
  \node[anchor=west] at (7.66,6.03) {\rev{Failed runs}};

  \node[anchor=east] at (-0.28,4) {\autospec};
  \node[anchor=east] at (-0.28,3) {\specgen};
  \node[anchor=east] at (-0.28,2) {Pure LLM};
  \node[anchor=east] at (-0.28,1) {\toolname\ w/o refine};
  \node[anchor=east] at (-0.28,0) {\textbf{\toolname}};

  \node[anchor=south] at (1.9,4.55) {Time (s)};
  \draw[figaxis] (0,-0.45) -- (3.8,-0.45);
  \node[anchor=north] at (0,-0.55) {0};
  \node[anchor=north] at (3.8,-0.55) {\rev{1200}};
  \node[anchor=south] at (6.8,4.55) {Tokens};
  \draw[figaxis] (4.9,-0.45) -- (8.7,-0.45);
  \node[anchor=north] at (4.9,-0.55) {0};
  \node[anchor=north] at (8.7,-0.55) {120k};
  \node[anchor=south] at (11.7,4.55) {API calls};
  \draw[figaxis] (9.8,-0.45) -- (13.6,-0.45);
  \node[anchor=north] at (9.8,-0.55) {0};
  \node[anchor=north] at (13.6,-0.55) {15};

  % Solid segment = successful cost / all attempts; hatched = failed cost / all attempts.
  \coststack{0.000}{4}{1.307}{2.952}{figgrayfill}{figgrayframe}{932.1}
  \coststack{0.000}{3}{0.445}{1.189}{figgrayfill}{figgrayframe}{375.5}
  \coststack{0.000}{2}{0.008}{0.039}{figgrayfill}{figgrayframe}{12.3}
  \coststack{0.000}{1}{0.051}{0.129}{figloopfill}{figloopstroke}{40.8}
  \coststack{0.000}{0}{0.856}{1.008}{figcontractfill}{figcontractstroke}{318.2}

  \coststack{4.900}{4}{5.789}{6.297}{figgrayfill}{figgrayframe}{44.1k}
  \coststack{4.900}{3}{4.966}{5.039}{figgrayfill}{figgrayframe}{4.4k}
  \coststack{4.900}{2}{4.913}{4.952}{figgrayfill}{figgrayframe}{1.6k}
  \coststack{4.900}{1}{4.955}{5.129}{figloopfill}{figloopstroke}{7.2k}
  \coststack{4.900}{0}{7.166}{7.553}{figcontractfill}{figcontractstroke}{83.8k}

  \coststack{9.800}{4}{11.859}{13.072}{figgrayfill}{figgrayframe}{12.9}
  \coststack{9.800}{3}{10.022}{10.269}{figgrayfill}{figgrayframe}{1.9}
  \coststack{9.800}{2}{9.867}{10.053}{figgrayfill}{figgrayframe}{1.0}
  \coststack{9.800}{1}{9.927}{10.304}{figloopfill}{figloopstroke}{2.0}
  \coststack{9.800}{0}{11.625}{11.896}{figcontractfill}{figcontractstroke}{8.3}
\end{tikzpicture}

\caption{\rev{Mean cost per attempted program on all 400 \texttt{gpt-5.4-mini} cases. Each bar stacks valid-run (solid) and failed-run (hatched) contributions, each divided by 400.}}
\label{fig:rq7-cost}
\end{figure*}

\rev{Each solid (hatched) bar segment sums valid-run (failed-run) costs over
all 400 attempts. Failed runs contribute 48.1 of \toolname's 318.2 seconds
per attempt, versus 519.4/932.1 for \autospec\ and 235.0/375.5 for
\specgen. Time per obtained valid specification is 338.6, 1,462.1, and
758.6 seconds, respectively. \toolname's 83.8k tokens per attempt include
repeated refinement context without cache discounts.}

\begin{finding}
\textbf{Finding (RQ7).} \rev{\toolname's added framework complexity mainly appears as higher token use; its wall-clock time and API-call count are much closer to the baselines than the raw token totals alone suggest.}
\end{finding}
